% Generated draft: requires mathematical and source review.
\documentclass[11pt]{article}
\usepackage[margin=25mm]{geometry}
\usepackage{fontspec}
\setmainfont{Latin Modern Roman}
\usepackage{amsmath,amssymb,mathtools,unicode-math}
\setmathfont{Latin Modern Math}
\usepackage{enumitem,needspace,xurl,hyperref}
\setlength{\parindent}{0pt}
\setlength{\parskip}{5pt}
\setlength{\emergencystretch}{3em}
\newcommand{\N}{\mathbb{N}}
\newcommand{\Z}{\mathbb{Z}}
\newcommand{\Q}{\mathbb{Q}}
\newcommand{\R}{\mathbb{R}}
\newcommand{\C}{\mathbb{C}}
\newcommand{\F}{\mathbb{F}}
\newcommand{\A}{\mathbb{A}}
\newcommand{\PP}{\mathbb{P}}
\newcommand{\E}{\mathbb{E}}
\newcommand{\eps}{\varepsilon}
\newcommand{\dd}{\,\mathrm{d}}
\begin{document}
\section*{UnsolvedMath ID: 30006508}
\textbf{Source problem code:} catalogue entry
\subsection{Definitions and mathematical statement}
A convex body in \(\R^n\) is a compact convex set with nonempty interior, where \(n\) is a positive integer. No symmetry assumption is imposed in the reported nonsymmetric setting.

The supplied source-verification field attributes the following estimate to Theorem 2 of the cited report, on report page 2943: there exists a universal constant \(C>0\), independent of the dimension and the body, such that every convex body \(K\subset\R^n\) in the report's isotropic position satisfies
\[
M(K)\leq \frac{C\log n}{\sqrt n}.
\]
This is a transcription of the theorem as reconstructed by the supplied verification field; the original extracted formula is corrupted. The complete supplied record does not define \(M(K)\), specify the normalization of isotropic position, or give the admissible range of dimensions. Consequently, the display is a reported estimate with unresolved definitions and hypotheses, rather than a fully specified mathematical statement. In particular, \(n=1\) makes its right-hand side zero, so a dimension restriction or another source qualification must be checked rather than silently supplied.

The surrounding source text describes obtaining a similar result in the nonsymmetric case as a formerly open problem and states that Pisier's argument fails there; the antecedent result and the relevant argument are not supplied. According to the supplied verification, the report then presents the displayed theorem and does not pose removal of the logarithm as a new problem. The separate supplied literature assessment describes the exact optimal estimate and remaining logarithmic factors as unresolved, but gives no precise conjecture, normalization, or quantified optimization target. Thus no distinct open problem can be formulated precisely from this record.
\subsection{Short English statement}
The supplied verification attributes a universal bound \(M(K)\leq C\log(n)/\sqrt n\) for nonsymmetric convex bodies in isotropic position to a theorem in the report. The definition of \(M\), the isotropic normalization, and the dimension range are missing. The record supplies no precise remaining open question; its partially solved label and literature assessment require independent review.
\subsection{Sources}
\begin{itemize}
\item Local UnsolvedMath registry, \texttt{data/unsolvedmath/problems.json}, permanent ID \(30006508\), source code \texttt{catalogue entry}, title: \emph{Reported Mean-Norm Bound for Nonsymmetric Convex Bodies}. The supplied record labels the status \texttt{partially\_solved} and the statement status \texttt{unrecoverable}; these are registry classifications, not independently verified mathematical conclusions. No proposer is identified; the supplied proposed year is 2025.
\item \emph{Functional Inequalities: Geometric Calculus meets Stochastic Analysis} (2026), Oberwolfach Reports, pp. 2917--2994, DOI \texttt{10.4171/owr/2025/54}: \url{https://doi.org/10.4171/owr/2025/54}. The supplied extraction refers to report page 2943 and Theorem 2. The report is supplied as the authoritative reference, but its contents have not been independently checked for this draft. The citation year 2026 and report identifier 2025/54 are retained as supplied.
\item Bizeul, \emph{Distances between non-symmetric convex bodies}: \url{https://arxiv.org/abs/2510.20511}. This is a supplied literature reference; no theorem from it is independently asserted here.
\item Supplied literature source: \emph{Oberwolfach Report 2025/54}, \url{https://doi.org/10.4171/owr/2025/54}.
\item Literature assessment checked at \texttt{2026-08-22}: the supplied curation describes an almost optimal general bound for the nonsymmetric analogue, with the exact optimal estimate and remaining logarithmic factors unresolved. This is supplied dated curation derived from the OWR record, not an independently verified literature assessment. The supplied background explicitly requests independent verification of time-sensitive status and cited claims.
\end{itemize}
\end{document}
